\begin{abstract}
Knowledge updating---modifying a Large Language Model's (LLM) internal facts post-training---is a critical capability for keeping models current. However, the side effects of such updates on the model's broader belief network and uncertainty estimation remain poorly understood. 
In this work, we investigate the "ripple effects" of injecting conflicting knowledge using a rigorously constructed, strictly bijective knowledge graph (20,000 nodes, Wikidata-verified). 
Contrary to the assumption that knowledge editing is local or that errors decay rapidly with distance, we report three destabilizing phenomena: 
(1) \textbf{Global Confidence Inflation}: The update process decouples confidence from accuracy, causing the model to become overconfident not only in its hallucinations but also in unrelated correct facts; 
(2) \textbf{The Damage Plateau}: The impact of injected errors does not vanish exponentially but maintains a steady error rate ($\sim$11\%) deep into the graph structure (up to 5 hops); 
(3) \textbf{Semantic Destabilization}: Instead of cleanly overwriting a fact, the injection increases the semantic entropy of outputs, leading to chaotic rather than targeted hallucinations. 
Our findings suggest that current metrics for knowledge editing are insufficient, as high "editing success" may mask profound structural degradation and calibrated "fake confidence."
\end{abstract}